\section{Development process}
\label{sec:process}

Development was organised as a sequence of numbered experiments (\#1--\#45) recorded in an
experiment log (\code{training/EXPERIMENTS.md}). Each entry records the hypothesis, the variants,
the measured numbers and the decision. Rejected ideas are logged too, so they are not retried.
\Cref{tab:versions} lists the versions that were adopted. Scores are
$(\text{wins} + \tfrac12\text{draws})/\text{games}$; ``Colab speed'' means the search budget was
halved to simulate the slower grading CPU (\cref{sec:protocol}).

\begin{table}[!htbp]
\centering\small
\caption{Adopted versions. ``Hard'' is the Hard AI of wrongway.app (\cref{sec:opponents}).
H2H = head-to-head against the previous adopted version.}
\label{tab:versions}
\begin{tabularx}{\linewidth}{@{}l X X@{}}
\toprule
Version (exp.) & Main change & Key measurement \\
\midrule
v0 (\#1) & greedy along the arena distance map, heuristic walls & loses to the sample \code{wall\_builder} \\
v1 (\#3) & alpha-beta, iterative deepening, tuple cells, depth $\le 4$ & beats greedy/random; still loses to \code{wall\_builder} (spends walls early) \\
v2 (\#5--6) & integer cells, frozenset of blocked edges, transposition table (depth 5--6); 30 generations of weight evolution & 0.85 vs v1; only 0.17--0.23 vs Hard \\
v4 (\#15--16) & vulnerability feature $u$; defensive walls at the root & 0.52 vs Hard (120 games), 0.74 vs v2 \\
v5 (\#19--22) & bitboard BFS ($2.05\times$, same moves); partial iterations; move ordering ($1.3\times$); progress tie-break & H2H 0.60 (160 games, Colab speed); vs Hard standard start $0.35\to0.47$ \\
v6 (\#23, \#26--27) & lead/reserve term (keep walls when ahead); vulnerability over 8 steps; rush rule & H2H 0.69--0.70 (reserve); exploit bot ``funnel'' 0 wins in 80 \\
v7 (\#31--34) & symmetric evaluation + tempo; depth cap $6\to10$; incremental distances ($2.6\times$) & Hard standard start 0.75, random 0.76; H2H 0.625 (160 games) \\
v8 (\#37) & evaluation rewritten as 428-entry tables, refined by TD-leaf($\lambda$) & H2H 0.49 (200 games); Hard 0.76 (200) / 0.74 (240) \\
\textbf{v9 (\#44--45)} & \textbf{held walls valued $1.2\times$ higher in the wall table} & vs wall-hoarding bots $0.23\to0.50$ (40 games); Hard standard start 0.775 (80); H2H 0.455 (160) --- \textbf{submitted} \\
\bottomrule
\end{tabularx}
\end{table}

\subsection{The turning point: depth, not evaluation}
\label{sec:depth}
Until v6, most attempts to improve the evaluation or to search unevenly failed:
\begin{itemize}
  \item a wall quiescence search (\#7);
  \item late-move reductions (\#24);
  \item search extensions for heavy walls (\#30).
\end{itemize}
Two diagnostics explained why.

\paragraph{Depth is the bottleneck (\#29).}
Ignoring the time limit, the same evaluation searched to a fixed depth of 6, 7 and 8 plies scored
0.59, 0.55 and \textbf{0.81} against Hard from the standard start (80 games each). Eight plies are
enough to see four consecutive opponent walls. Better heuristics could not replace that.

\paragraph{The evaluation had an odd/even bias (\#31).}
The root value at depth $d$ minus the value at depth $d-1$ was $+9.7, -9.3, +9.5, -9.4$ for
$d = 3..6$ (74 positions). For comparison, one step of distance is worth about 1.38 points.
The cause was two distance weights applied relative to the side to move at the leaf: for the two
players' views of the same position, $E(p) + E(q) = -0.43\,(d_p + d_q) \approx -9$ instead of~0.
The agent's style therefore flipped with the parity of the depth it happened to reach. This
explained the back-and-forth shuffling seen in replays and the failure of every uneven-depth
technique (they compare odd and even branches).
\begin{itemize}
  \item \textbf{Fix (\#32).} Dropping the two weights and adding a tempo constant reduced the bias
        to below 0.4.
  \item \textbf{Payoff (\#33).} With the bias gone, a deeper cap finally helped: Hard
        standard-start score rose from 0.56 to 0.70 with cap 8. Cap 10 was kept.
\end{itemize}
From then on, effort went into speed (\#34, $2.6\times$ with identical candidate sets).

\subsection{A standard-start weakness: undervalued walls (\#44--45)}
\label{sec:walls}
Two bots from our own team (\code{dtest2}, \code{Phat\_test2}) share 95\% of their code and hold
their walls for later. Other teams were known to follow the same idea, so we treated them as a
proxy for that style.
\begin{itemize}
  \item \textbf{The problem.} From the standard start, with both sides at simulated grading speed
        and games played one at a time, v8 lost all 20 games against them. As first player it
        replayed the same 74-move loss every time.
  \item \textbf{Diagnosis (\#44).} A deeper search of our own (depth 9) agreed with every move
        from move 37 on, so the game was lost earlier. v8 had spent 7 of its 10 walls by move 25,
        against 3 for the opponent, and the opponent's remaining walls then pushed our path from 7
        to 13 steps.
  \item \textbf{Two things ruled out.} A fixed depth of 7 or 8 still lost as first player, so an
        opening book computed by deep search would not help. The opponent also searched
        \emph{shallower} than us (7.15 vs 8.0 plies on average), so it did not win by depth.
  \item \textbf{Conclusion.} The tables learned by self-play value a wall in hand too low against
        an opponent that, unlike our own engine, keeps its walls.
  \item \textbf{Fix (\#45).} We scaled the held-wall part of the table,
        $W[w][d] \leftarrow W[0][d] + s\,(W[w][d]-W[0][d])$. Each value of $s$ was measured on
        standard-start games at ten different time budgets, to obtain different game lines.
\end{itemize}
The results for $s = 1.0 / 1.2 / 1.4$:
\begin{itemize}
  \item against the wall-hoarding pair: 0.23 / 0.50 / 0.65 (40 games each);
  \item against Hard from the standard start: 0.75 / 0.775 / 0.625;
  \item head-to-head against $s=1$: 0.455 / 0.395 (160 games).
\end{itemize}
$s=1.4$ only exploits one style, at a cost against Hard. $s=1.2$ fixes the weakness without a
measurable loss elsewhere (0.966 vs 0.969 against eight other bots), so it was adopted.

\subsection{Rejected ideas}
\Cref{tab:rejected} lists the main rejected ideas. A recurring lesson: optimising against a single
opponent, or training at a shallower depth than the agent plays at, gives weights that do not
transfer.

\begin{table}[!htbp]
\centering\small
\caption{Main rejected ideas (all measured with the protocol of \cref{sec:protocol}).}
\label{tab:rejected}
\begin{tabularx}{\linewidth}{@{}l X X@{}}
\toprule
Idea (exp.) & Result & Reason \\
\midrule
Value network, 1-ply (\#2) & 0--5\% vs heuristic bots & learned the spurious rule ``placing walls wins'' \\
Wall quiescence (\#7) & 0.51--0.62 vs v2 & costs depth \\
Repetition penalty (\#8) & 0.45 vs v2 & treats a symptom, not the cause \\
\code{pass} in the tree (\#12--13) & 0.26--0.30 vs Hard, 0.30--0.50 vs v2 & larger tree, no gain \\
Evolution against Hard only (\#14, \#17) & better vs Hard or vs alpha-beta bots, never both & over-fits one opponent; trained at depth 4 \\
Late-move reductions (\#24) & 0.35 vs previous & reduced moves were exactly the trap walls \\
Wall-pair trap search (\#25) & +0.14 vs Hard standard start, $-0.10$ random & not better everywhere \\
Heavy-wall extensions (\#30, \#33) & 0.33--0.38 vs previous & uneven depth \\
Re-evolution after v7 (\#35) & 1 of 14 generations accepted; H2H 0.47 & weights already at a local optimum \\
TD at depth 2--3 (\#36) & H2H 0.34 & depth mismatch between training and play \\
TD ``stance'' tables, +202 params (\#38) & H2H 0.44 (120 games) & no gain, 4--5\% slower \\
AlphaZero-style policy/value nets (\#39--41) & best 0.49 vs MSSV (60 games) & too little self-play data in a pure-Python budget \\
Shared wall-effect cache (\#42) & identical moves, same speed & cache hit rate only 29\% \\
Keep TT and caches across moves (\#43) & average depth 6.36 vs 6.59 & the last iteration dominates the time \\
Opening book from deep search (\#44) & depth 7--8 still loses the same opening & problem is evaluation, not depth \\
Wall scale $s=1.4$ (\#45) & 0.625 vs Hard, H2H 0.395 & exploits one style only \\
\bottomrule
\end{tabularx}
\end{table}

\subsection{AI-assisted variant pipeline}
\label{sec:pipeline}
The course allows AI tools. We used an AI coding assistant (Claude Code by Anthropic) throughout
the project. The team set the goals and made every adoption decision. The assistant wrote most
variants, ran the measurements and kept the log. Every change to the submitted file was reviewed,
and the team can explain it. The pipeline was:
\begin{enumerate}
  \item \textbf{Hypothesis from failures.} Losing games were saved as replays. The position
        before the mistake was rebuilt, and the root scores of every candidate were printed at
        each depth. Examples: the 50-turn shuffling (\#22) and the odd/even bias (\#31).
  \item \textbf{Variant generation.} The assistant implemented each hypothesis as a copy of the
        current agent behind a flag or parameter, for example \code{PROGRESS\_BONUS}
        $\in\{0.001, 0.2\}$ or \code{VULN\_STEPS} $\in\{4, 8, \infty\}$. Parameter sweeps usually
        produced 2--5 files at once. Optimisation-only changes were also checked for identical
        moves with a fixed hash seed.
  \item \textbf{Static gates.}
        \begin{itemize}
          \item Unit tests: the simulated rules must match the arena's moves, legal walls and
                distances on more than 500 random positions; self-play must stay legal; the
                agent must take an immediate win.
          \item A submission checker: both sides against the three sample bots, zero
                timeouts/errors, and the longest move time reported.
        \end{itemize}
  \item \textbf{Measurement} with the gauntlet of \cref{sec:protocol}: 40--240 games per
        opponent, at simulated grading speed.
  \item \textbf{Selection criteria.} A variant was adopted only if
        \begin{enumerate}[label=(\alph*)]
          \item it was not worse than the current champion head-to-head beyond noise (about one
                standard deviation, \cref{sec:protocol});
          \item it was not worse against Hard, from both the standard start and random openings;
          \item it still won the standard-start games against every sparring bot;
          \item it was clearly better somewhere.
        \end{enumerate}
        When two versions were equal, the simpler one was kept (\#38).
\end{enumerate}

\paragraph{Scale.}
\begin{itemize}
  \item 45 logged experiments containing about 100 hand-designed variants.
  \item Five evolution runs: 86 generations of 6 mutants, i.e.\ about 500 weight vectors, each
        judged on several dozen games.
  \item Three TD-leaf runs (14\,400, 4\,000 and 4\,000 self-play games).
  \item Three AlphaZero-style runs.
\end{itemize}
In total the agent played tens of thousands of games during development.
